\subsection{Local cohomology}

In this number, we consider a local Noetherian ring A. Recall (VIII, § 3, n° 3, lemma 2) that the ideals of definition of A are the ideals of A distinct from A containing a power of \(\mathfrak{m}_{\mathrm{A}}\) , or equivalently the ideals \(\mathfrak{a} \subset \mathfrak{m}_{\mathrm{A}}\) such that \(\mathrm{A}/\mathfrak{a}\) is of finite length. We will denote \(\mathcal{D}\) the set of defining ideals of \(\mathrm{A}\) , equipped with the order relation opposite to inclusion; it is filtered on the right.

Let M be an A-module. Associate with every ideal of definition a of A the graded A-module \(\operatorname{Ext}_{\mathrm{A}}(\mathrm{A}/\mathfrak{a},\mathrm{M})\) ; if a and b are ideals of definition with \(a \subset b\) , let us denote by \(p_{ab}: A/a \to A/b\) the canonical map, and consider the A-linear map \(\operatorname{Ext}(p_{ab},1_{\mathrm{M}}): \operatorname{Ext}_{\mathrm{A}}(\mathrm{A}/\mathfrak{b},\mathrm{M}) \longrightarrow \operatorname{Ext}_{\mathrm{A}}(\mathrm{A}/\mathfrak{a},\mathrm{M})\) . We thus obtain an inductive system of graded A-modules and graded A-linear maps of degree 0 relative to the ordered set D. We call the local cohomology module of M, and denote \(\mathrm{H}_{\mathrm{A}}(\mathrm{M})\) , the graded A-module \(\lim \operatorname{Ext}_{\mathrm{A}}(\mathrm{A}/\mathfrak{a},\mathrm{M})\) .

The ideals \(\mathfrak{m}_{\mathrm{A}}^{n}\) for \(n \geqslant 1\) form a cofinal subset of \(\mathcal{D}\) ; we therefore have a canonical isomorphism of graded modules \(\varinjlim_{n} \operatorname{Ext}_{\mathrm{A}}(\mathrm{A}/\mathfrak{m}_{\mathrm{A}}^{n}, M) \longrightarrow \mathrm{H}_{\mathrm{A}}(\mathrm{M})\) . Consequently, every element of \(\mathrm{H}_{\mathrm{A}}(\mathrm{M})\) is annihilated by a power of \(\mathfrak{m}_{\mathrm{A}}\) .

Remark 1.--- *Let X be the topological space* \(\operatorname{Spec}(A)\) , \(\mathcal{O}_{\mathrm{X}}\) the structural sheaf of rings and \(\widetilde{\mathbf{M}}\) the \(\mathcal{O}_{\mathrm{X}}\) -module associated with M. The graded A-module \(\mathrm{H}_{\mathrm{A}}(\mathrm{M})\) is identified with the module \(\mathrm{H}_{\{\mathfrak{m}_{\mathrm{A}}\}}(\mathrm{X},\widetilde{\mathrm{M}})\) of cohomology with support in the closed point \(\mathfrak{m}_{\mathrm{A}}\) of \(X\) .

For every homomorphism \(f: \mathrm{M} \to \mathrm{N}\) of A-modules, the maps \(\operatorname{Ext}_{\mathrm{A}}(1_{\mathrm{A}/\mathfrak{a}}, f): \operatorname{Ext}_{\mathrm{A}}(\mathrm{A}/\mathfrak{a}, \mathrm{M}) \longrightarrow \operatorname{Ext}_{\mathrm{A}}(\mathrm{A}/\mathfrak{a}, \mathrm{N})\) form an inductive system of graded linear maps. Passing to the inductive limit, one obtains a graded homomorphism \(\mathrm{H}_{\mathrm{A}}(f): \mathrm{H}_{\mathrm{A}}(\mathrm{M}) \to \mathrm{H}_{\mathrm{A}}(\mathrm{N})\) . For every sequence \(\mathrm{M} \xrightarrow{f} \mathrm{N} \xrightarrow{g} \mathrm{P}\) of A-modules and homomorphisms, one has \(\mathrm{H}_{\mathrm{A}}(g \circ f) = \mathrm{H}_{\mathrm{A}}(g) \circ \mathrm{H}_{\mathrm{A}}(f)\) . Let

\[
0 \longrightarrow \mathrm{M} \stackrel {f} {\longrightarrow} \mathrm{N} \stackrel {g} {\longrightarrow} \mathrm{P} \longrightarrow 0\tag{ℰ}
\]

an exact sequence of A-modules. By A, X, p.90, prop. 8, the connecting homomorphisms of the extension modules \(\operatorname{Ext}_{\mathrm{A}}(\mathrm{A}/\mathfrak{a},\mathrm{P}) \longrightarrow \operatorname{Ext}_{\mathrm{A}}(\mathrm{A}/\mathfrak{a},\mathrm{M})\) form an inductive system of A-linear maps, graded of (ascending) degree +1. Passing to the inductive limit, one deduces an A-homomorphism \(\partial(\mathcal{E}): \mathrm{H}_{\mathrm{A}}(\mathrm{P}) \to \mathrm{H}_{\mathrm{A}}(\mathrm{M})\) , graded of degree +1, which makes exact the sequence of homomorphisms

\[
\dots \longrightarrow \mathrm{H} _ {\mathrm{A}} ^ {n - 1} (\mathrm{P}) \xrightarrow {\partial^ {n - 1} (\mathcal {E})} \mathrm{H} _ {\mathrm{A}} ^ {n} (\mathrm{M}) \xrightarrow {\mathrm{H} _ {\mathrm{A}} ^ {n} (f)} \mathrm{H} _ {\mathrm{A}} ^ {n} (\mathrm{N}) \xrightarrow {\mathrm{H} _ {\mathrm{A}} ^ {n} (g)} \mathrm{H} _ {\mathrm{A}} ^ {n} (\mathrm{P}) \xrightarrow {\partial^ {n} (\mathcal {E})} \mathrm{H} _ {\mathrm{A}} ^ {n + 1} (\mathrm{M}) \longrightarrow \dots
\]

Let M be an A-module. For every ideal a of A, the A-module \(\mathrm{Hom}_{\mathrm{A}}(\mathrm{A}/\mathfrak{a},\mathrm{M})\) is canonically identified with the submodule of M consisting of the elements annihilated by a. Thus \(\mathrm{H}_{\mathrm{A}}^{0}(\mathrm{M})\) identifies with the submodule of M consisting of the elements m that are annihilated by a power of \(\mathfrak{m}_{\mathrm{A}}\) , that is, such that \(\mathrm{long}_{\mathrm{A}}(\mathrm{Am}) < +\infty\) . In particular, we have \(\mathrm{H}_{\mathrm{A}}^{0}(\mathrm{M}) = \mathrm{M}\) when M is Artinian.

Examples.--- 1) If M is injective, the A-module \(\mathrm{H}_{\mathrm{A}}^{i}(\mathrm{M})\) is zero for \(i > 0\) and injective for \(i = 0\) ( \(\S 8, \mathrm{n}^{\circ}2\) , lemma 1, c)).

\begin{enumerate}
\def\labelenumi{\arabic{enumi})}
\setcounter{enumi}{1}
\item
  If \(\mathrm{H}_{\mathrm{A}}^{0}(\mathrm{M})=\mathrm{M}\) (for example if M is Artinian), \(\mathrm{H}_{\mathrm{A}}^{i}(\mathrm{M})\) is zero for i > 0. Let (I, e) be an injective envelope of M. The submodule \(\mathrm{H}_{\mathrm{A}}^{0}(\mathrm{I})\) of I is injective (example 1) and contains \(e(\mathrm{M})\) , therefore is equal to I. Set N to be the cokernel of e and consider the exact sequence \(0\rightarrow M\xrightarrow{e}\mathrm{I}\xrightarrow{p}\mathrm{N}\rightarrow0\) . Since \(\mathrm{I}=\mathrm{H}_{\mathrm{A}}^{0}(\mathrm{I})\) , one has \(\mathrm{N}=\mathrm{H}_{\mathrm{A}}^{0}(\mathrm{N})\) and the homomorphism \(\mathrm{H}_{\mathrm{A}}^{0}(p)\) is surjective. Since \(\mathrm{H}_{\mathrm{A}}^{i}(\mathrm{I})\) is zero for i > 0 (example 1), \(\mathrm{H}_{\mathrm{A}}^{1}(\mathrm{M})\) is zero, and \(\mathrm{H}_{\mathrm{A}}^{i}(\mathrm{M})\) is isomorphic to \(\mathrm{H}_{\mathrm{A}}^{i-1}(\mathrm{N})\) for i > 1; we conclude by reasoning by induction on the integer i.
\item
  Let \(\Omega\) be a dualizing A-module. For \(i \neq \dim(A)\) , one has \(\operatorname{Ext}_{A}^{i}(A/a, \Omega) = 0\) for every ideal of definition \(a\) of A ( \(\S 8\) , \(n^{\circ} 5\) , example 3), whence \(\mathrm{H}_{\mathrm{A}}^{i}(\Omega) = 0\) ; for \(i = \dim(A)\) , the A-module \(\mathrm{H}_{\mathrm{A}}^{i}(\Omega)\) , which is isomorphic to \(\varinjlim \operatorname{Ext}_{\mathrm{A}}^{i}(\mathrm{A}/\mathfrak{m}_{\mathrm{A}}^{n}, \Omega)\) , is a Matlis A-module (loc. cit., example 6).
\item
  Let A be a Noetherian local integral domain; denote by K its field of fractions, and suppose A ≠ K. K is an injective A-module (A, X, p.18, example 1), so that the module \(\mathrm{H}_{\mathrm{A}}(K)\) is zero (example 1). From the exact sequence 0 → A → K → K/A → 0, we derive for every i an isomorphism \(\mathrm{H}_{\mathrm{A}}^{i}(K/A)\) → \(\mathrm{H}_{\mathrm{A}}^{i+1}(A)\).
\end{enumerate}

More generally, for every torsion-free A-module M and every integer i, we deduce from the exact sequence

\[
0 \rightarrow \mathrm{M} \rightarrow \mathrm{K} \otimes_ {\mathrm{A}} \mathrm{M} \rightarrow (\mathrm{K/A}) \otimes_ {\mathrm{A}} \mathrm{M} \rightarrow 0
\]

an isomorphism \(\mathrm{H}_{\mathrm{A}}^{i}((\mathrm{K}/\mathrm{A})\otimes_{\mathrm{A}}\mathrm{M})\to \mathrm{H}_{\mathrm{A}}^{i + 1}(\mathrm{M})\)

\begin{enumerate}
\def\labelenumi{\arabic{enumi})}
\setcounter{enumi}{4}
\tightlist
\item
  Keep the hypotheses of the preceding example and suppose moreover that dim(A) = 1. Let N be a torsion A-module; since every nonzero ideal of A distinct from A is an ideal of definition (VIII, § 1, n° 3, prop. 6, e)), we have \(\mathrm{H}_{\mathrm{A}}^{0}(N) = N\), and consequently \(\mathrm{H}_{\mathrm{A}}^{i}(N) = 0\) for i > 0 (example 2).
\end{enumerate}

Let M be an A-module; denote by T(M) its torsion submodule. Consider the long exact sequence of local cohomology associated with the exact sequence

\[
0 \to \mathrm{T} (\mathrm{M}) \to \mathrm{M} \to \mathrm{M} / \mathrm{T} (\mathrm{M}) \to 0;
\]

taking the preceding into account, we deduce from it canonical isomorphisms T(M) → \(\mathrm{H}_{\mathrm{A}}^{0}(M)\) and \(\mathrm{H}_{\mathrm{A}}^{1}(M)\) → \(\mathrm{H}_{\mathrm{A}}^{1}(M/T(M))\). Since the canonical homomorphism \((K/A) \otimes_{A} M\) → \((K/A) \otimes_{A} (M/T(M))\) is bijective, we finally obtain canonical isomorphisms

\[
\mathrm{H} _ {\mathrm{A}} ^ {0} (\mathrm{M}) \rightarrow \mathrm{T} (\mathrm{M}), \quad \mathrm{H} _ {\mathrm{A}} ^ {1} (\mathrm{M}) \rightarrow (\mathrm{K} / \mathrm{A}) \otimes_ {\mathrm{A}} \mathrm{M}.
\]

PROPOSITION 1. Let A be a Noetherian local ring and M a finitely generated A-module.

\begin{enumerate}
\def\labelenumi{\alph{enumi})}
\item
  The A-module \(\mathrm{H}_{\mathrm{A}}(\mathrm{M})\) is Artinian, and zero in degree \(>\dim (\mathrm{M})\) .
\item
  Set \(p = \mathrm{prof}_{\mathrm{A}}(\mathrm{M})\) . One has \(\mathrm{H}_{\mathrm{A}}^{i}(\mathrm{M}) = 0\) for \(i < p\) , and \(\mathrm{H}_{\mathrm{A}}^{p}(\mathrm{M}) \neq 0\) if \(M\) is nonzero.
\end{enumerate}

Let us prove a) by reasoning by induction on \(\dim(M)\) . The case \(\dim(M) \leqslant 0\) follows from example 2 above. Suppose \(\dim(M) > 0\) and first take \(M\) of the form \(A/p\) , where \(p\) is a prime ideal of \(A\) distinct from \(\mathfrak{m}_A\) . Let \(x\) be an element of \(\mathfrak{m}_A - p\) ; we have an exact sequence \(0 \to M \xrightarrow{x_M} M \longrightarrow M/xM \to 0\) , with \(\dim(M/xM) = \dim(M) - 1\) . From this we deduce an exact sequence of local cohomology

\[
\mathrm{H} _ {\mathrm{A}} ^ {i - 1} (\mathrm{M} / x \mathrm{M}) \longrightarrow \mathrm{H} _ {\mathrm{A}} ^ {i} (\mathrm{M}) \stackrel {x} {\longrightarrow} \mathrm{H} _ {\mathrm{A}} ^ {i} (\mathrm{M}).
\]

Every element of \(\mathrm{H}_{\mathrm{A}}^{i}(\mathrm{M})\) is annihilated by a power of \(\mathfrak{m}_{\mathrm{A}}\) ; to prove that this module is Artinian, it therefore suffices to prove that the socle of \(\mathrm{H}_{\mathrm{A}}^{i}(\mathrm{M})\) is of finite dimension over \(\kappa_{A}\) (§ 8, n° 3, lemma 3). By the induction hypothesis, the kernel N of the homothety of ratio x in \(\mathrm{H}_{\mathrm{A}}^{i}(\mathrm{M})\) is Artinian; since x belongs to \(\mathfrak{m}_{\mathrm{A}}\) , the socle of \(\mathrm{H}_{\mathrm{A}}^{i}(\mathrm{M})\) identifies with that of N, hence is of finite dimension. If i > dim(M), we have \(\mathrm{H}_{\mathrm{A}}^{i-1}(\mathrm{M}/x\mathrm{M}) = 0\) by the induction hypothesis, so that the homothety of ratio x is injective in \(\mathrm{H}_{\mathrm{A}}^{i}(\mathrm{M})\) ; since every element of \(\mathrm{H}_{\mathrm{A}}^{i}(\mathrm{M})\) is annihilated by a power of x, we deduce \(\mathrm{H}_{\mathrm{A}}^{i}(\mathrm{M}) = 0\) , whence a) in the case under consideration.

Let us pass to the general case. The A-module M admits a composition series \((\mathrm{M}_{j})_{0\leqslant j\leqslant n}\) such that each quotient \(M_{j}/M_{j+1}\) is isomorphic to \(A/p_{j}\) , where \(p_{j}\) is a prime ideal of A (IV, § 1, n° 4, th. 1). Let us prove by induction on n that M satisfies a). The case n=0 is trivial. The exact sequence \(0\to M_{1}\to M\to A/p_{0}\to0\) provides an exact sequence of local cohomology

\[
\mathrm{H} _ {\mathrm{A}} ^ {i} \left(\mathrm{M} _ {1}\right) \longrightarrow \mathrm{H} _ {\mathrm{A}} ^ {i} (\mathrm{M}) \longrightarrow \mathrm{H} _ {\mathrm{A}} ^ {i} \left(\mathrm{A} / \mathfrak {p} _ {0}\right);
\]

the A-module \(\mathrm{H}_{\mathrm{A}}^{i}(\mathrm{M}_{1})\) is Artinian by the induction hypothesis, and the same holds for \(\mathrm{H}_{\mathrm{A}}^{i}(\mathrm{A}/\mathfrak{p}_{0})\) by the cases already treated; consequently \(\mathrm{H}_{\mathrm{A}}^{i}(\mathrm{M})\) is Artinian. If \(i > \dim(\mathrm{M})\) the modules \(M_{1}\) and \(A/p_{0}\) are of dimension < i; the modules \(\mathrm{H}_{\mathrm{A}}^{i}(\mathrm{M}_{1})\) and \(\mathrm{H}_{\mathrm{A}}^{i}(\mathrm{A}/\mathfrak{p}_{0})\) are therefore zero by the induction hypothesis and the cases already treated, which implies \(\mathrm{H}_{\mathrm{A}}^{i}(\mathrm{M}) = 0\) .

Suppose M is nonzero, and prove b) by induction on the integer \(p = \text{prof}(M)\) . The case \(p = 0\) follows from the definition of depth. Suppose \(p > 0\) and let us choose an element \(x\) of \(\mathfrak{m}_{\mathrm{A}}\) such that the homothety \(x_{\mathrm{M}}\) is injective. As above, we obtain an exact sequence of local cohomology.

\[
\mathrm{H} _ {\mathrm{A}} ^ {i - 1} (\mathrm{M} / x \mathrm{M}) \longrightarrow \mathrm{H} _ {\mathrm{A}} ^ {i} (\mathrm{M}) \xrightarrow {x} \mathrm{H} _ {\mathrm{A}} ^ {i} (\mathrm{M}).
\]

We have \(\operatorname{prof}(M/xM) = \operatorname{prof}(M) - 1\) ( \(\S 1, n^{\circ} 4\) prop. 7), whence \(H_{A}^{i-1}(M/xM) = 0\) for \(i < p\) by the induction hypothesis, which implies as above \(H_{A}^{i}(M) = 0\) . In particular \(H_{A}^{p-1}(M)\) is zero, so that the homomorphism \(H_{A}^{p-1}(M/xM) \longrightarrow H_{A}^{p}(M)\) is injective; thus \(H_{A}^{p}(M)\) is nonzero by the induction hypothesis.

One can show that the module \(H_{A}^{\dim(M)}(M)\) is nonzero when M is nonzero (exerc. 4; cf. no. 3, cor. of th. 2).

COROLLARY. Let M be a nonzero finitely generated Macaulay A-module. The A-module \(\mathrm{H}_{\mathrm{A}}^{i}(\mathrm{M})\) is zero for \(i \neq \dim(\mathrm{M})\) and nonzero for \(i = \dim(\mathrm{M})\) .

Remark 2.--- For every ideal of definition a of A, the A-module \(\operatorname{Ext}_{\mathbf{A}}(\mathbf{A}/\mathfrak{a},\mathbf{M})\) is annihilated by a, and A/a is identified with \(\widehat{\mathbf{A}}/\mathfrak{a}\widehat{\mathbf{A}}\) ; therefore, the graded A-module \(\operatorname{Ext}_{\mathbf{A}}(\mathbf{A}/\mathfrak{a},\mathbf{M})\) is identified with \(\widehat{\mathbf{A}}\otimes_{\mathbf{A}}\operatorname{Ext}_{\mathbf{A}}(\mathbf{A}/\mathfrak{a},\mathbf{M})\) , and hence also with \(\operatorname{Ext}_{\widehat{\mathbf{A}}}(\widehat{\mathbf{A}}/\mathfrak{a}\widehat{\mathbf{A}},\widehat{\mathrm{A}}\otimes_{\mathbf{A}}\mathbf{M})\) (A, X, p.111, prop. 10). The set of ideals \(\mathfrak{a}\widehat{\mathbf{A}}\) , for \(\mathfrak{a}\in\mathcal{D}\) , contains the powers of \(\mathfrak{m}_{\widehat{\mathbf{A}}}\) , hence is cofinal in the set of defining ideals of \(\widehat{\mathbf{A}}\) ; we therefore deduce from the preceding an isomorphism of graded A-modules

\[
\mathrm{H} _ {\mathrm{A}} (\mathrm{M}) \longrightarrow \mathrm{H} _ {\widehat {\mathrm{A}}} (\widehat {\mathrm{A}} \otimes_ {\mathrm{A}} \mathrm{M}).
\]

If the A-module M is of finite type, the A-module \(\widehat{\mathrm{A}}\otimes_{\mathrm{A}}\mathrm{M}\) is identified with the completion \(\widehat{\mathrm{M}}\) of M (III, § 3, n° 4, th. 3), and we have an isomorphism \(\mathrm{H}_{\mathrm{A}}(\mathrm{M})\to \mathrm{H}_{\widehat{\mathrm{A}}}(\widehat{\mathrm{M}})\) , graded of degree 0.

